% Einheit 1 – Klammern auflösen (terme.md Einheit 3) – Nr. 10–16
\einheitenkopf[e1][1 Klammern auflösen]{Einheit 1 von 5 · Klammern auflösen}
\zweigzeile{Hier lernst du, Klammern in Termen aufzulösen · neu oder schon bekannt – je nach Buch · keine P10-Aufgabe · baut auf: Glieder zusammenfassen, Rechnen mit Vorzeichen}
\setcounter{aufgabe}{9}

\begin{aufgabe}{Ich erkenne, welches Vorzeichen zu welchem Glied gehört.}
\anweisung{Schreibe die Glieder mit ihrem Vorzeichen auf.}
\begin{geruest}
\gz{a}{$6x + 2y - 5$}{\leerfeld \quad \leerfeld \quad \leerfeld}
\gz{b}{$3x - 4 + y$}{\leerfeld \quad \leerfeld \quad \leerfeld}
\gz{c}{$-x + 7y - 2$}{\leerfeld \quad \leerfeld \quad \leerfeld}
\gz{d}{$8 - 3x - y$}{\leerfeld \quad \leerfeld \quad \leerfeld}
\gz{e}{$-2x - 5y + 9$}{\leerfeld \quad \leerfeld \quad \leerfeld}
\end{geruest}
\end{aufgabe}

\begin{aufgabe}{Ich erkenne, was vor einer Klammer steht.}
\anweisung{Kreuze an, was direkt vor der Klammer steht.}
\begin{teile}
\teil $5 + (x - 2)$ \hfill \kreuz{Plus} \kreuz{Minus} \kreuz{Zahl}
\teil $9x - (4 + x)$ \hfill \kreuz{Plus} \kreuz{Minus} \kreuz{Zahl}
\teil $3 \cdot (x + 1)$ \hfill \kreuz{Plus} \kreuz{Minus} \kreuz{Zahl}
\teil $2x - 4 \cdot (x - 3)$ \hfill \kreuz{Plus} \kreuz{Minus} \kreuz{Zahl}
\teil $7 - (2x - 1)$ \hfill \kreuz{Plus} \kreuz{Minus} \kreuz{Zahl}
\end{teile}
\end{aufgabe}

\begin{aufgabe}{Ich kann Klammern auflösen (neu in diesem Jahr).}
\anweisung{Löse die Klammer auf.}
\rechenplatz{2}
\begin{gleichungsraster}[1]
\gl{x + (y + 4)} & \gl{3x + (y - 2)} \\
\gl{5 + (2x - y)} & \gl{3 \cdot (x + 5)} \\
\gl{2 \cdot (y + 6)} & \gl{5 \cdot (x + 2)} \\
\gl{4 \cdot (y + 7)} & \gl{4x - (y + 2)} \\
\gl{x - (3y - z + 4)} & \gl{-3 \cdot (x - 4)} \\
\gl{x \cdot (x + 6)} & \\
\end{gleichungsraster}
\end{aufgabe}

\begin{aufgabe}{Ich kann Klammern auflösen und den Term zusammenfassen.}
\anweisung{Löse die Klammern auf und fasse zusammen.}
\rechenplatz{2}
\begin{gleichungsraster}[2]
\gl{3x + (2x + 4)} & \gl{6 + 2 \cdot (x + 1)} \\
\gl{5y + 3 \cdot (y - 2)} & \gl{7y - (2y + 5)} \\
\gl{10 - 2 \cdot (x - 4)} & \gl{2 \cdot (x + 3) + 4 \cdot (x - 1)} \\
\gl{3x \cdot (x - 2) - 2 \cdot (x^2 - 4x + 5)} & \\
\end{gleichungsraster}
\end{aufgabe}

\begin{aufgabe}{Ich finde den Fehler beim Auflösen einer Klammer.}
\anweisung{Finde den Fehler, benenne ihn und korrigiere die Rechnung.}
\begin{teile}
\teil Tom rechnet: \rechnung{9x - (4x - 6) &= 9x - 4x - 6 \\ &= 5x - 6}
\teil Mia rechnet: \rechnung{12 - 3 \cdot (x - 2) &= 12 - 3x - 6 \\ &= 6 - 3x}
\end{teile}
\end{aufgabe}

\begin{aufgabe}{Ich kann begründen, ob zwei Terme gleichwertig sind.}
\anweisung{Entscheide, ob die beiden Terme gleichwertig sind, und begründe.}
\begin{teile}
\teil $2 \cdot (x + 5)$ \ und \ $2x + 5$
\teil $10 - (x - 3)$ \ und \ $13 - x$
\teil $-(4 - x)$ \ und \ $x - 4$
\end{teile}
\end{aufgabe}

\begin{aufgabe}{Ich kann eine Sachaufgabe mit Klammern lösen.}
Ein rechteckiger Garten ist $(x + 6)$\,m lang und $x$\,m breit. Er bekommt rundherum einen Zaun, nur für das Tor bleiben $3$\,m frei.

\rechteck[seiten={x+6,x,,},punkte=]{4.5}{2.5}

\begin{teile}
\teil Stelle einen Term für die Länge des Zauns auf und vereinfache ihn.
\teil Der Garten ist $8$\,m breit. Wie viele Meter Zaun braucht man?
\end{teile}
\end{aufgabe}
